\documentclass[../../main.tex]{subfiles}
\begin{document}

{\bf [解答]}

表記の簡潔さのため$s = \sin x, c = \cos x$，また明らかな時は$f(y)=f$などと引数を略記する．
与えられた関数方程式
\begin{align}
  f(x)=\dfrac{a}{2\pi}\int_0^{2\pi}\sin(x+y)f(y)dy+
  \dfrac{b}{2\pi}\int_0^{2\pi}\cos(x-y)f(y)dy+\sin x+\cos x \label{eq:1}
\end{align}
について考える．式中に現れる積分について$x,y$を分離すると
\begin{align}
  \begin{cases}
    \int_0^{2\pi} \sin(x+y) f(y) \, dy = s \int_0^{2\pi} f\cos y \, dy + c \int_0^{2\pi} f\sin y \, dy
    \int_0^{2\pi} \cos(x-y) f(y) \, dy = c \int_0^{2\pi} f\cos y \, dy + s \int_0^{2\pi} f\sin y \, dy
  \end{cases} \label{eq:2}
\end{align}
である．ここで簡単のため
\begin{align}
  \begin{cases}
    P = \int_0^{2\pi} f(y)\cos y \, dy \\
    Q = \int_0^{2\pi} f(y)\sin y \, dy \\
  \end{cases} \label{eq:3}
\end{align}
とおくと\cref{eq:2}に代入して
\begin{align}
  \begin{cases}
    \int_0^{2\pi} \sin(x+y) f(y) \, dy = Ps+Qc \\
    \int_0^{2\pi} \cos(x-y) f(y) \, dy = Pc+Qs
  \end{cases} \label{eq:4}
\end{align}
である．\cref{eq:1}に代入して
\begin{align}
  f(x)
  &= \left[ \dfrac{a}{2\pi} P + \dfrac{b}{2\pi} Q + 1 \right] s \\
  &\quad + \left[ \dfrac{a}{2\pi} Q + \dfrac{b}{2\pi} P + 1 \right] c  \label{eq:5}
\end{align}
である．
さらに左の括弧内を $A_{s}$，右の括弧内を $A_{c}$ として:
\begin{align}
  \begin{cases}
    A_{s} = \dfrac{a}{2\pi} P + \dfrac{b}{2\pi} Q + 1  \\
    A_{c} = \dfrac{a}{2\pi} Q + \dfrac{b}{2\pi} P + 1
  \end{cases} \label{eq:6}
\end{align}
\cref{eq:5}に代入すると
\begin{align}
  f &= A_{s} s + A_{c} c \label{eq:7}
\end{align}
である．\cref{eq:3}の右辺に代入すると
\begin{align}
  P
  &= \int_0^{2\pi} (A_{s} sc + A_{c} c^2) \, dx \\
  &= -\dfrac{1}{4} A_{s} \left[ \cos 2x \right]_0^{2\pi} + \dfrac{1}{2} A_{c} \left[ x + \dfrac{1}{2}\sin 2x \right]_0^{2\pi} \\
  &= A_{c} \pi \label{eq:8}
\end{align}
および
\begin{align}
  Q
  &= \int_0^{2\pi} (A_{s} s^2 + A_{c} cs) \, dx \\
  &= \dfrac{1}{2} A_{s} \left[ x - \dfrac{1}{2}\sin 2x \right]_0^{2\pi} - \dfrac{1}{4} A_{c} \left[ \cos 2x \right]_0^{2\pi} \\
  &= A \pi \label{eq:9}
\end{align}
となる．\cref{eq:8,eq:9}を$A,B$の定義式\cref{eq:6}に代入すると
\begin{align}
  &
  \begin{cases}
    A_{s} = \dfrac{a}{2\pi} A_{c} \pi + \dfrac{b}{2\pi} A_{s} \pi + 1 = \dfrac{1}{2} a A_{c} + \dfrac{1}{2} b A_{s} + 1 \\
    A_{c} = \dfrac{a}{2\pi} A_{s} \pi + \dfrac{b}{2\pi} A_{c} \pi + 1 = \dfrac{1}{2} a A_{s} + \dfrac{1}{2} b A_{c} + 1
  \end{cases} \\
  \therefore
  &
  \begin{cases}
    (b-2)A_{s} + aA_{c} + 2 = 0 \\
    aA_{s} + (b-2)A_{c} + 2 = 0
  \end{cases} \label{eq:10}
\end{align}
である．

さて，\cref{eq:7}から $(A_{s},A_{c})$ が一意に定まることが $f$ が一意に定まることの必要十分条件である．そこで\cref{eq:10}で$A_{s},A_{c}$が一意に定まるような$a,b$の条件を求める．
まず$(a,b)=(0,2)$の時は不適だからこれ以外で考える．すると\cref{eq:10}の二式は$A_{s}A_{c}$平面における直線を表し，解$(A_{s},A_{c})$は二直線の交点である．よってこれが一意に定まるのは二直線が平行ではない時で，条件は
\begin{align}
  (b-2)^2 &\neq a^2 \quad \therefore b-2 \neq \pm a \label{eq:11}
\end{align}
である．このもとで\cref{eq:10}を解くと
\begin{align}
  A_{s} = A_{c} = \dfrac{2}{2-(a+b)} \label{eq:12}
\end{align}
であり，\cref{eq:7}に代入して求める$f(x)$は
\begin{align}
  f(x) &= \dfrac{2}{2-(a+b)} (\sin x + \cos x)
\end{align}
である．

{\bf [別解]}

\cref{eq:10}からの議論は以下のようにしても良いだろう．
\cref{eq:10}の和差をとって
\begin{align}
  \begin{cases}
    (b-2+a)(A_{s}+A_{c}) = -4 \\
    (b-2-a)(A_{s}-A_{c}) = 0
  \end{cases} \label{eq:13}
\end{align}
である．

\subparagraph{$1^\circ \; b-2-a \neq 0$の時}

\cref{eq:13}の第二式から $A_{s}=A_{c}$ である．
第一式より
\begin{align}
  A_{s} = A_{c} = \dfrac{2}{2-(a+b)}
\end{align}
であり，この時 $f$ は一意に定まる．

\subparagraph{$2^\circ \; k-a = 0$の時}

\cref{eq:13}の第二式は任意の $(A,B)$ で成立する．
第一式から
\begin{align}
  A_{c} &= -A_{s} - \dfrac{2}{k+a}
\end{align}
だが，これは複数の$(A_{s},A_{c})$で成立するから$f$ は一意に定まらない．
\casesend

以上から，求める条件は $b-2 \neq \pm a$ つまり $f(x) = \dfrac{2}{2-(a+b)}(\sin x + \cos x)$ である．

% TODO
% 元の手書き回答にはあったけど，ここの議論これって必要か？
%
% $(A_{s},A_{c}) = (0,0)$ の時は\cref{eq:10}は成立せず不適だから，$A_{s}, A_{c}$ のうち少なくとも一方は $0$ でない．したがって

\end{document}
